\scanpage{62}

\diagram{\begin{tikzcd}
\Delta/X\arrow[r,"\text{fibration}"]&\Delta
\end{tikzcd}}

\deleted{
\diagram{\begin{tikzcd}
C\arrow[r]&B\\
C^\circ\arrow[d,"F"']&\\
\deleted{\unsure{(\mathrm{Ab})}}\inserted{\mathcal E}&
\end{tikzcd}}
\diagram{\begin{tikzcd}
C_b&i^*(x')\arrow[d]&b\arrow[d,"i"]\\
C_{b'}\arrow[u,"i^*"']&x'&b'
\end{tikzcd}}

\margin{(\unsure{adjoint})}
\diagram{\begin{tikzcd}[column sep=small]
(b\in B)\arrow[r,mapsto]&
{\displaystyle\unsure{\sum|\prod}_{x\in C_b}\unsure{F|R}(x)}
\arrow[r]&F(\unsure{i^*x'})\\
(b'\in B)\arrow[r,mapsto]&
{\displaystyle\sum_{x'\in C_{b'}}\unsure{F|R}(x')}
\arrow[u]\arrow[r]&
F(\deleted{\unsure{i^*}}\unsure{x'})\arrow[u]
\end{tikzcd}}
\deleted{$\unsure{\downarrow}$}

\diagram{\begin{tikzcd}[column sep=small]
(b\in B)\arrow[r,mapsto]&
{\displaystyle\prod_{x\in C_b}\unsure{F|R}(x)}
\arrow[d,"\unsure{?}"']\arrow[r]&
F(i^*x')\arrow[d]\\
(b'\in B)\arrow[r,mapsto]&
{\displaystyle\prod_{x'\in C_{b'}}F(x')}
\arrow[r]&F(x').
\end{tikzcd}}
\margin{$\unsure{f^*}$}
}

\margin{\unsure{?}}

Si
\diagram{\begin{tikzcd}
(\Delta/X)^\circ\arrow[r,"A_*"]&
\deleted{(\mathrm{Ab})}\inserted{\mathcal E\text{ (cat. avec sommes directes)}}
\end{tikzcd}}
est un faisceau sur $X$ (i.e. \deleted{\unsure{faisceau}}\inserted{syst.} de coeff. \uline{\unsure{contravariant}} \unsure{par r. à} un simplex variable sur $X$)
on lui associe
\[
C_*(X,A_*)=i_!(A):(\Delta\deleted{\unsure{/X}})^\circ\mathbin{-}\mathcal E
\]
(où $i:\widehat{\Delta/X}\longrightarrow\widehat{\Delta}$ est le morphisme d'induction des topos)

Si
\diagram{\begin{tikzcd}
\Delta/X\arrow[r,"A^*"]&
\deleted{(\mathrm{Ab})}\inserted{\mathcal E\text{ (cat. av. produits)}}
\end{tikzcd}}
est un ``cofaisceau'' sur $X$ (i.e. syst. de coeff. \uuline{cov.} \unsure{par r. à} un simplex variable de $X$) on lui associe $C^*(X,A^*)$
\[
\deleted{\unsure{co}}:(\Delta\deleted{\unsure{{}^\circ}})\mathbin{-}\mathcal E.
\]

On passe de l'un à l'autre en remplaçant $\mathcal E$ par $\mathcal E^\circ$; c'est compatible avec composition avec foncteur $\mathcal E\longrightarrow\mathcal E'$ qui commute aux sommes resp. produits.

\margin{61}
